% 3.4 — Les principes qui gouvernent les décisions
\section{Les principes qui gouvernent les décisions}
\label{sec:principes}

Chaque principe ci-dessous est présenté de la même façon : l'énoncé, la commande qui casse
l'approche naïve, la règle retenue et le test qui la fige. Les sorties de terminal sont
réelles (commit \commitref{}, Linux, sans filet d'exécution sauf mention contraire) ; seules
les lignes sans intérêt sont coupées et les chemins abrégés en \cmd{\textasciitilde/projet}.

% ------------------------------------------------------------------
\subsection{Analyser et simuler plutôt que reconnaître des motifs}
\label{sec:pourquoi-moteur}

\paragraph{Le constat.} La bêta et la version 1.7.0 lisaient les commandes avec un
découpeur de texte et des expressions régulières. L'audit et la revue menés au début de la
période documentée ont trouvé, encore et encore, des défauts de la même forme : \emph{deux
analyseurs répondant différemment à la même question} --- le découpeur de lignes, le
normaliseur de mot de commande, le masqueur de citations, le lecteur de \textit{heredocs},
le découpeur d'opérandes. Deux de ces défauts avaient été introduits par des corrections
des autres (\fichier{tests/test\_command\_model.py}).

\paragraph{Ce qui casse.} Le shell n'est pas un texte mais une grammaire. Un commentaire
neutralise ce qui le suit (\cmd{ls \# ; rm -rf x} ne supprime rien) ; une substitution
\cmd{\$(\ldots)} exécute une commande cachée dans un argument (\cmd{git commit -m "\$(rm
y)"}) ; une enveloppe déplace le vrai mot de commande (\cmd{sudo -u x rm},
\cmd{bash -o pipefail -c "rm app.db"}, \cmd{env -S 'rm -rf x'}) ; un \cmd{case} ou une
fonction contiennent des commandes que le découpage par \cmd{;} ne voit pas ; une citation
ANSI-C (\cmd{\$'\textbackslash x72m'}) écrit \cmd{rm} sans les lettres \cmd{r} et \cmd{m}.

\paragraph{La règle retenue.} Une seule analyse de la ligne (\fichier{model.py}), lue par
tout le reste. Avant de remplacer le découpeur, elle a tourné en ombre : sur toutes les
commandes de la suite de tests (801, dans les deux dialectes), elle a rendu la même décision
que le découpeur ; sur un jeu étiqueté à la main de ce que bash exécute réellement, elle
avait raison dans 22 cas où le découpeur avait tort, et tort dans aucun cas où il avait
raison. Ces deux chiffres justifient le remplacement plutôt que la cohabitation.

\paragraph{Et la simulation ?} Une structure correcte ne suffit pas
(section~\ref{sec:statique}) : \cmd{mv a.txt b.txt} n'est destructeur que si \cmd{b.txt}
existe. Le moteur rejoue donc la ligne sur une vue virtuelle du système de fichiers et suit
l'origine des contenus : dans \cmd{mv a b \&\& rm b}, ce qui est perdu est le contenu de
\cmd{a}. C'est exactement le raisonnement qui manquait à l'agent dans l'incident Gemini CLI.

% ------------------------------------------------------------------
\subsection{L'inconnu n'est pas sûr : « non couvert » et « non compris »}
\label{sec:inconnu}

\paragraph{Deux ignorances différentes.} Le moteur distingue :
\begin{itemize}
  \item \textbf{non couvert} (\textit{uncovered}) : la ligne est lue, mais un outil n'a pas
    de règle. On sait \emph{ce qui} est appelé, avec quels arguments ;
  \item \textbf{non compris} (\textit{not understood}) : une partie de la ligne ne peut pas
    être lue avant l'exécution --- une variable affectée ailleurs, une substitution dont la
    sortie n'est pas prévisible, un mot de commande construit.
\end{itemize}

\paragraph{La règle retenue (parties M2.3a et M2.3b, arbitrages M1--M5 et Rb1--Rb4).}
\begin{enumerate}
  \item Une ligne \textbf{non comprise} va au filet d'exécution s'il est disponible, sinon à
    un point de contrôle s'ils sont activés, sinon elle est \textbf{escaladée}. Elle
    n'obtient jamais ALLOW.
  \item Un \textbf{exécuteur de code} (\cmd{python}, \cmd{node}, \cmd{make}, \cmd{npm},
    \cmd{./script.sh}, \cmd{source}\ldots) lance du code que \demo{} ne lit pas : il va au
    filet, sinon au point de contrôle, sinon il est autorisé \emph{avec une note qui dit
    qu'aucune capture n'a été faite} (arbitrage M1). Il garde le filet même si on lui passe
    un fichier hors du projet (V21).
  \item Deux cas d'exécuteur sont plus risqués et \textbf{escaladent sans filet} (supervisés
    avec, arbitrage Rb2) : du \textbf{code en ligne qui lance d'autres programmes}
    (\cmd{python -c "\ldots os.system(\ldots)"}, \cmd{node -e "\ldots child\_process\ldots"})
    et les \textbf{installations de paquets} (\cmd{pip install}, \cmd{npm install}\ldots), qui
    exécutent les scripts d'installation de chaque paquet.
  \item Une ligne qui atteint quelque chose au-delà des fichiers locaux sans être lue
    (cloud, conteneurs, services, une option inconnue sur un outil distant, \cmd{kubectl
    exec}, \cmd{heroku run}\ldots) est escaladée : ni le filet ni le point de contrôle ne la
    voient. Le code récupéré du réseau puis exécuté (\cmd{curl \ldots{} | bash}) garde son
    refus.
  \item Un mot sans règle qui n'est pas un exécuteur de code n'a pas d'effet (arbitrage Rb1).
\end{enumerate}

\begin{terminal}[ligne non comprise, pas de filet]
$ demo_cli check --mode enforce 'rm $INCONNU'
demo_cli 1.7.0  [x] BLOCKED  ·  ESCALATE
  matched rule        not_understood
Decision
  No recovery path for a not_understood on 'unknown'; cannot auto-recover.
  Human input required.
$ demo_cli check --mode enforce 'python3 outil.py'
demo_cli 1.7.0  [+] SAFE  ·  ALLOW
Decision
  Allowed: it runs code demo_cli does not read (python3); no runtime net was
  available, so nothing that code changes is captured or watched.
$ demo_cli check --mode enforce 'pip install requests'
demo_cli 1.7.0  [x] BLOCKED  ·  ESCALATE
Decision
  No recovery path for a unread_code on 'unknown'; cannot auto-recover. Human input
  required. It installs packages, which run their own install scripts (pip), and no
  runtime net was available to capture what that changes.
\end{terminal}

\begin{pourquoi}[Pourquoi l'exécuteur sans filet est-il autorisé ?]
Escalader chaque \cmd{python script.py}, \cmd{make} ou \cmd{npm test} rendrait l'outil
inutilisable : c'est l'essentiel du travail d'un agent. L'arbitrage M1 accepte ce cas \emph{à
visage découvert} : la raison de la décision dit qu'aucune capture n'a été faite. Les deux cas
où l'exécuteur lance à coup sûr du code étranger --- un programme démarré depuis du code en
ligne, une installation de paquets --- escaladent en l'absence de filet (Rb2). Là où le filet
existe (Linux avec Claude Code, Windows avec le montage), ces lignes sont SUPERVISED. Un
outil inconnu qui n'est pas un exécuteur de code (« non couvert ») est, lui, traité comme
sans effet (Rb1) ; la table de règles s'étend famille par famille pour réduire ce cas, et un
test (\fichier{test\_m22f.py}) vérifie qu'aucune ligne du portail sur laquelle une ancienne
règle se déclenchait n'est restée hors du moteur sans être marquée non comprise.
\end{pourquoi}

\paragraph{Lire plutôt qu'abandonner.} « Non compris » doit rester rare, sinon tout escalade.
Le moteur lit donc ce qui peut l'être sans rien exécuter : la sortie de \cmd{\$(cat
liste.txt)}, \cmd{echo}, \cmd{printf}, \cmd{pwd}, \cmd{basename}, \cmd{dirname}, sous des
conditions strictes (fichier régulier du projet, UTF-8, au plus 64~Kio, non modifié plus tôt
sur la ligne ; partie I1b). La partie M2.3b a ajouté ce que l'ancien chemin savait capturer :
une valeur connue sur la ligne lue comme l'option qu'elle est (\cmd{X=-rf; rm \$X src}), un
script littéral envoyé à un shell par un tube (\cmd{echo 'rm -rf build' | sh}),
\cmd{sqlite3 \$DB} et \cmd{sqlite3 app.db "\$(cat q.sql)"}, les messages de commit, et
\cmd{npx}, \cmd{bunx} ou \cmd{pnpm dlx} qui lancent un outil connu (\cmd{npx prisma migrate
reset} est lu comme \cmd{prisma}) :

\begin{terminal}[sortie de substitution lue par le moteur]
$ demo_cli check --mode enforce 'rm $(cat list.txt)'
demo_cli 1.7.0  [+] SAFE  ·  REVERSIBLE
  target              ~/projet/a.txt
Recovery
  snapshot            a.txt.20261009-175242.6d6ca4f4.bak
Decision
  Recovery point captured before the action; reversible.
\end{terminal}

% ------------------------------------------------------------------
\subsection{Jamais un REVERSIBLE non mérité}
\label{sec:jamais-reversible}

\paragraph{L'énoncé.} REVERSIBLE est une promesse : « \cmd{demo\_cli undo} remettra tout en
place ». Une capture partielle, ou une copie qui ne contient pas ce que la commande détruit,
rendrait cette promesse fausse, ce qui est pire qu'une escalade : la personne cesse de se
méfier.

\paragraph{Les règles qui en découlent.}
\begin{itemize}
  \item \textbf{Tout ou rien} : les captures d'une commande forment un groupe ; si l'une échoue,
    les autres sont invalidées et la commande est escaladée (partie S2,
    \cmd{recovery.snapshot\_set}).
  \item \textbf{Pas de point de contrôle pour ce qu'il ne contient pas} : une copie du projet
    ne restaure ni un fichier hors du projet (V17), ni une base distante, ni un historique
    git poussé de force.
  \item \textbf{Hors projet, rien n'est promis} : une destruction connue hors du projet est
    escaladée, même à côté d'une partie non comprise.
  \item \textbf{Le décideur ne capture pas} : \cmd{decide.py} ne reçoit qu'un booléen « une
    capture de la cible réelle a réussi », fourni par l'orchestrateur ; il n'existe aucune
    cible par défaut sur laquelle se rabattre.
\end{itemize}

% ------------------------------------------------------------------
\subsection{Pendant la migration, la réponse la plus stricte gagne}
\label{sec:plus-strict}

\paragraph{Le problème.} Le moteur a remplacé l'ancien chemin \emph{progressivement}, famille
par famille. Pendant ce temps, deux lecteurs coexistent : le moteur et les anciennes règles.
Que faire quand ils divergent ?

\paragraph{La règle retenue (parties B0 et M2.3a).} L'ancien chemin peut garder une
\emph{capture} ou une \emph{escalade}, jamais son ALLOW. Exemple réel, au moment de M2.3a : le
moteur marquait \cmd{X=-rf; rm \$X build} comme non compris (une valeur qui commence par un
tiret) ; l'ancienne règle, elle, savait capturer \cmd{build}. La réponse la plus stricte ---
capturer --- l'emportait. À l'inverse, quand l'ancienne règle aurait autorisé une ligne que le
moteur ne lit pas, c'est la règle M3 qui s'applique (filet, point de contrôle, escalade) : la
vulnérabilité V19 était précisément un \cmd{psql -c "\$(cat q.sql)"} autorisé alors que
\cmd{q.sql} contenait \cmd{DROP TABLE users;}. Il est désormais escaladé.

\paragraph{La fin de la migration (M2.3b).} La partie M2.3b a fait lire au moteur tout ce que
l'ancien chemin capturait encore ; depuis, \textbf{l'ancien chemin ne décide plus aucune ligne
POSIX}. La même ligne est maintenant capturée par le moteur lui-même :

\begin{terminal}[la ligne est lue par le moteur (M2.3b)]
$ demo_cli check --mode enforce 'X=-rf; rm $X build'
demo_cli 1.7.0  [+] SAFE  ·  REVERSIBLE
  target              ~/projet/build
Recovery
  snapshot            build.20261009-175242.30656068.snapdir
Decision
  Recovery point captured before the action; reversible.
\end{terminal}

La règle « le plus strict gagne » ne vaut plus que là où l'ancien chemin sert encore : le texte
\cmd{cmd.exe}, une ligne POSIX qui nomme une cmdlet PowerShell (M5), et le cas où le moteur
échoue (B0). Une construction PowerShell que C3 ne lit pas
n'y retourne pas : elle est non comprise, et suit le même chemin qu'une ligne POSIX non lue
--- filet, sinon point de contrôle, sinon escalade.

% ------------------------------------------------------------------
\subsection{Les fichiers de \texorpdfstring{\demo}{demo\_cli} : le plan de contrôle}
\label{sec:plan-controle}

\paragraph{L'énoncé.} Un agent qui peut modifier la configuration de son garde-fou n'a pas de
garde-fou. Le \emph{plan de contrôle} regroupe : la configuration \fichier{.demo\_cli.toml},
l'espace de travail \fichier{.demo\_cli/} (registre, réserve de récupération), les
configurations des \textit{hooks} dans le projet et dans le répertoire personnel, le script du
garde du shell, \fichier{\textasciitilde/.bashrc} quand il l'active, le lanceur
\cmd{demo\_cli}, et sous Windows les fichiers de la tâche d'ouverture de session.

\paragraph{La règle.} Toute modification de ces chemins --- y compris une \emph{création} ou un
ajout en fin de fichier --- est refusée (ESCALATE, jamais approuvable sur place), à toutes les
portes : avant l'exécution par le moteur et par les anciennes règles, dans la porte fichier,
et pendant l'exécution par le garde d'appels système, qui annule l'appel.

\begin{terminal}[un agent tente de passer l'outil en mode observation]
$ demo_cli check --mode enforce 'echo mode = "shadow" > .demo_cli.toml'
demo_cli 1.7.0  [x] BLOCKED  ·  ESCALATE
  matched rule        control_plane
Decision
  This changes demo_cli's own files, or a directory holding them (.demo_cli.toml):
  its config, hook settings, ledger or recovery store. An agent may not [...]
\end{terminal}

% ------------------------------------------------------------------
\subsection{Demander ou refuser}
\label{sec:ask-deny}

\paragraph{Le problème.} Une escalade arrête l'agent. Faut-il alors \emph{demander} à la
personne, ou \emph{refuser} ? Demander tout use l'attention (section~\ref{sec:permissions}) ;
refuser tout rend l'outil inutilisable pour des opérations ordinaires (supprimer un volume
Docker de test, pousser une branche).

\paragraph{La règle retenue (partie ASK).} L'escalade reste une escalade --- aucune décision du
portail ne change --- mais chaque surface porte un \emph{canal} dans
\fichier{policy.ESCALATION\_CHANNEL} :
\begin{itemize}
  \item \textbf{ask} quand le changement est réel mais ordinaire, et que la commande exacte suffit
    à la personne pour savoir à quoi elle consent (base distante, historique git, bucket,
    conteneur, déploiement, service\ldots) ;
  \item \textbf{deny} quand tout part sans retour réaliste (infrastructure détruite, vidage d'un
    entrepôt de données, formatage de disque, paiement, compte utilisateur supprimé, données
    de récupération du système) ; seuls un jeton d'approbation structurel, que l'agent ne
    peut pas forger, ou une exécution hors de l'agent le permettent.
\end{itemize}
\textbf{Une surface sans ligne dans la table est refusée} : une surface ajoutée plus tard ne
devient jamais « demandable » par accident. Enfin, \emph{demander} n'a de sens que si
quelqu'un voit la question : le \textit{hook} Claude Code ne pose la question que dans les
modes où l'interface l'affiche (\cmd{default}, \cmd{acceptEdits}, \cmd{plan}, \cmd{auto}) et
refuse dans les autres --- une question que personne ne voit ne doit jamais devenir « c'est
passé ». Codex, qui ne sait pas demander, refuse.

% ------------------------------------------------------------------
\subsection{Les faits sur les outils sont vérifiés}
\label{sec:faits-verifies}

Une règle fausse est une décision fausse. Chaque grammaire d'options est tirée de la source
de l'outil lui-même --- son \cmd{--help}, sa page de manuel, son code source, ses données
publiées (l'index des commandes d'\cmd{aws}, l'arbre de commandes de \cmd{gcloud}, la table
de commandes de \cmd{redis}) --- et vérifiée sur la version courante. Les comportements
douteux sont testés sur des fichiers jetables : que \cmd{unzip -Pxo} lit \cmd{xo} comme mot de
passe et non comme l'option d'écrasement (I1c, vérifié sur UnZip 6.0), que \cmd{helm 4}
change la sémantique de \cmd{rollback --dry-run}, que \cmd{tofu} accepte \cmd{-chdir DIR}
sans \cmd{=} quand \cmd{terraform} le refuse. Un test confronte chaque grammaire au
\cmd{--help} de l'outil installé. Une règle dont la grammaire vient seulement de la
documentation porte \cmd{verified=False} : c'est le cas de \nbnonverifiees{} règles sur
\nbregles{}.
